% lint-skip: 未完成ステータス（figure/table caption 未対応，人手チェック時に解消予定）
\documentclass[../../main.tex]{subfiles}
\begin{document}

\begin{flushright}
  \footnotesize\textit{【自動文字起こし・要確認】}
\end{flushright}

[解] $\dfrac{1}{3}$ の確率でAとBが，$\dfrac{2}{3}$ の確率でAとCが札を交換する

\paragraph{(1)}

\begin{figure}[htbp]
  \centering
  \begin{tikzpicture}[x=1.2cm,y=1cm,>=stealth,
    every node/.style={font=\small}]
    \node[gray!75] at (0,3.55) {$n$ 回目};
    \node[gray!75] at (3,3.55) {$n+1$ 回目};
    \begin{scope}[yshift=2.3cm]
      \node[draw=blue!75!black,rounded corners,fill=blue!7] (a0) at (0,0) {$a_n$};
      \node[draw=teal!70!black,rounded corners,fill=teal!7] (ab) at (3,0.55) {$b_{n+1}$};
      \node[draw=orange!85!black,rounded corners,fill=orange!7] (ac) at (3,-0.55) {$c_{n+1}$};
      \draw[->,blue!75!black,thick] (a0.east)--(ab.west) node[midway,above] {$1/3$};
      \draw[->,blue!75!black,thick] (a0.east)--(ac.west) node[midway,below] {$2/3$};
      \node[left] at (-0.42,0) {A が赤};
    \end{scope}
    \begin{scope}
      \node[draw=teal!70!black,rounded corners,fill=teal!7] (b0) at (0,0) {$b_n$};
      \node[draw=blue!75!black,rounded corners,fill=blue!7] (ba) at (3,0.55) {$a_{n+1}$};
      \node[draw=teal!70!black,rounded corners,fill=teal!7] (bb) at (3,-0.55) {$b_{n+1}$};
      \draw[->,teal!70!black,thick] (b0.east)--(ba.west) node[midway,above] {$1/3$};
      \draw[->,teal!70!black,thick] (b0.east)--(bb.west) node[midway,below] {$2/3$};
      \node[left] at (-0.42,0) {B が赤};
    \end{scope}
    \begin{scope}[yshift=-2.3cm]
      \node[draw=orange!85!black,rounded corners,fill=orange!7] (c0) at (0,0) {$c_n$};
      \node[draw=blue!75!black,rounded corners,fill=blue!7] (ca) at (3,0.55) {$a_{n+1}$};
      \node[draw=orange!85!black,rounded corners,fill=orange!7] (cc) at (3,-0.55) {$c_{n+1}$};
      \draw[->,orange!85!black,thick] (c0.east)--(ca.west) node[midway,above] {$2/3$};
      \draw[->,orange!85!black,thick] (c0.east)--(cc.west) node[midway,below] {$1/3$};
      \node[left] at (-0.42,0) {C が赤};
    \end{scope}
  \end{tikzpicture}
  \caption{赤い札を持つ人の，1回の交換による状態遷移}
\end{figure}

上図から
\begin{align}
  a_{n+1} &= \dfrac{1}{3} b_n + \dfrac{2}{3} c_n \\
  b_{n+1} &= \dfrac{1}{3} a_n + \dfrac{2}{3} b_n \\
  c_{n+1} &= \dfrac{2}{3} a_n + \dfrac{1}{3} c_n
\end{align}

\paragraph{(2)}

$a_n + b_n + c_n = 1$ から $c_n$ を消去し(1)に代入して
\begin{equation}
  a_{n+1} = \dfrac{1}{3}(-2a_n - b_n + 2) \tag{①}
\end{equation}
\begin{equation}
  b_{n+1} = \dfrac{1}{3}a_n + \dfrac{2}{3}b_n \tag{②}
\end{equation}
①から $b_n = -3a_{n+1} - 2a_n + 2$ を②に代入
$$-3a_{n+2} - 2a_{n+1} + 2 = \dfrac{1}{3}a_n + \dfrac{2}{3}(-3a_{n+1} - 2a_n + 2)$$
$$9a_{n+2} = 3a_n + 2$$
$$a_{n+2} - \dfrac{1}{3} = \dfrac{1}{3} \left(a_n - \dfrac{1}{3}\right)$$

これと $a_0 = 1, a_1 = 0$ から，等比数列の公式から
$$a_{2k} = \left(\dfrac{1}{3}\right)^k \dfrac{2}{3} + \dfrac{1}{3}$$
$$a_{2k+1} = \left(\dfrac{1}{3}\right)^k \left(-\dfrac{1}{3}\right) + \dfrac{1}{3}$$
だから
$$a_n =
\begin{cases} \dfrac{1}{3}\left[2\left(\dfrac{1}{3}\right)^{\frac{n}{2}} + 1\right] & (n \in \text{even}) \\ \dfrac{1}{3}\left[-\left(\dfrac{1}{3}\right)^{\frac{n-1}{2}} + 1\right] & (n \in \text{odd})
\end{cases}$$

\end{document}
